\Needspace{9\baselineskip}
\originalchapter{补充题与两次测验}
附加题原件有2.111J / 18.435J / ESD.79联列, PS3、4、7、8、9明确注明Fall 2004, 尽管被2003站点资源索引托管, 不冒称2003当年作业. 其余附加页未印年, 只保留托管版本. 以下按实际文件分组, 各编号要求均可对照原件; 教材4.26只有题号而留作缺题面; 5.6所对应公开文件已给Fourier/相移加法任务, 本卷独立重述和解答, 不复刊原书段落.
\originalsection{附加PS1、PS2: 可逆门与Pauli}
\source{Optional PS1, 第1--3题, 第2题含加分.}
\begin{exercise}(1) 用Fredkin及任意固定0、1工作位实现NOT、AND、fan-out, 证经典普适; (2) 用Toffoli建两个输入半加器, 加分含进位Z的全加器; (3) 验Pauli等价各轴$\pi$旋转、平方$I$、循环乘积$XY=\ii Z,YZ=\ii X,ZX=\ii Y$, 不同Pauli反对易.\end{exercise}
\begin{solution}(1) Fredkin控制a、目标0,1时输出目标为a,1-a, 同时得到复制基态与NOT. 控制a、目标b,0时末端为ab, 得AND并保留垃圾保证可逆; 配NOT与AND可实现一般经典布尔网络. (2) 工作位c=0,Toffoli(X,Y;c)给c=XY; 以X控制Y作CNOT得到和X$\oplus$Y. 全加器设新位S,C均0, S经三CNOT得X$\oplus$Y$\oplus$Z, C依次三个Toffoli(X,Y)、(X,Z)、(Y,Z)得XY$\oplus$XZ$\oplus$YZ, 是多数进位. (3) $R_i(\pi)=-\ii\sigma_i$, 整体相位无关; 直接矩阵乘积或$\sigma_i\sigma_j=\delta_{ij}I+\ii\epsilon_{ijk}\sigma_k$同时验证所有关系.
\end{solution}
\source{Optional PS2, 第1--6题.}
\begin{exercise}(1) 单位向量$\mathbf j$, 求$\sigma_j^2$, 本征值与归一本征态; (2) $\delta^2\ll1$时展开$\e^{\ii\delta Y/2}\e^{\ii\delta X/2}\e^{-\ii\delta Y/2}\e^{-\ii\delta X/2}$至二阶; (3) Bell正态在X乘积基的概率; (4) 对任意双体纯直积交换两端的矩阵; (5) $\bigl(\ket{0}_A\ket{1}_C+\ket{1}_A\ket{0}_C\bigr)\ket+_B/\sqrt2$在ABC顺序的三位计算基展开; (6) 简化${}_A\bra+\bigl(\ket{0}_A\ket{1}_B+\ket{1}_A\ket{0}_B\bigr)/\sqrt2$.\end{exercise}
\begin{solution}(1) 参数$\mathbf j=(\sin\theta\cos\phi,\sin\theta\sin\phi,\cos\theta)$, 态$\cos(\theta/2)0+\e^{\ii\phi}\sin(\theta/2)1$和$-\e^{-\ii\phi}\sin(\theta/2)0+\cos(\theta/2)1$, 值$\pm1$, 平方I. $\sigma_j$对应$\pi$旋转至全局相位, 与原题一致. (2) 群对易子$=I+[\ii\delta Y/2,\ii\delta X/2]+O(\delta^3)=I+\ii\delta^2Z/2+O(\delta^3)$. (3) Bell正态在X基为$(++ + --)/\sqrt2$, 同号各1/2, 异号0. (4) SWAP固定00,11、交换01,10. (5) 先重排寄存器至ABC, 得$(\ket{001}+\ket{011}+\ket{100}+\ket{110})/2$, 其余四项系数为零. (6) $\bra+0\rangle=\bra+1\rangle=1/\sqrt2$, 所以结果为$(\ket0_B+\ket1_B)/2=\ket+_B/\sqrt2$. 这是未归一的测量分支, 范数平方$1/2$, 归一后为$\ket+_B$.
\end{solution}

\originalsection{附加PS3、PS4: 关联与查询}
\source{Optional PS3, 第1--5题.}
\begin{exercise}(1) $U^{(j)}=(\sigma_j\otimes I+I\otimes\sigma_j)/2$, $V=\sum_j[U^{(j)}]^2$, 求四Bell态期望; (2) 求$[XX,YY]$, $[YY,ZZ]$及$[ZZ,XX]$; (3) 单态对一端取迹; (4) 验CNOT自伴且平方I; (5) qutrit R,Z相位与循环T移位的对易子和超密编码.\end{exercise}
\begin{solution}(1) $V=3I/2+(XX+YY+ZZ)/2=I+\mathrm{SWAP}$, 单态值0, 三个triplet Bell态均2. (2) 三者彼此对易, 因两个位置各反对易, 双负号抵消. (3) 单态约化$I/2$. (4) 计算基置换是10、11互换的对合, 矩阵实对称. (5) $RT=\omega TR$, 故$[R,T]=(\omega-1)TR$; $R^aT^b$的9个Bell像由迹正交得到9正交态, 每qutrit可传两trit, 资源须含预共享对.
\end{solution}
\source{Optional PS4, 第1--6、8题, 第7题商业教材引用.}
\begin{exercise}(1) 验$\sigma=Z\cos\theta+X\sin\theta$本征态$\cos(\theta/2)0+\sin(\theta/2)1$及正交态; (2) Bell正态在此两实基同时旋转后形式不变; (3) 一端Z、一端上述基,$\theta=\pi/3$, 求四联合概率; (4) GHZ的XXX、XYY、YXY、YYX期望; (5) 一一f及逆f分别有$C_1,C_2$门且产生工作垃圾, 构造$\ket x\ket0\to\ket{f(x)}\ket0$且成本$O(C_1+C_2+n)$; (6) CNOT在X乘积基矩阵; (8) 从bit-query构造phase-query.\end{exercise}
\begin{solution}(1) 乘矩阵得本征值$\pm1$. (2) 同一实正交R两端作用保持$\sum_i\ket{ii}$, 因$RR^T=I$. (3) 同号各$\cos^2(\theta/2)/2=3/8$, 异号各1/8. (4) $+1,-1,-1,-1$, 经典预定值乘积矛盾见主作业. (5) 先可逆求f, 复制正交f标签到新寄存器, 逆C1清垃圾; 再用逆f电路从f重新求x, 与旧x XOR消零并逆转逆f工作空间, 最后交换输出寄存器. 两电路各常数次, 复制/交换$O(n)$. 这只复制计算基标签, 不克隆未知态. (6) H两端共轭CNOT使控制目标互换, 矩阵固定00,10, 交换01,11, 此处基次序$++, +-, -+, --$. (8) 目标置$-$, 查询后相位$(-1)^{f(x)}$, 目标保持$-$可丢弃.
\end{solution}

\originalsection{附加PS5、PS6: 变换与多控制}
\source{Optional PS5, 第1--4题.}
\begin{exercise}(1) 对归一态
\[\ket\psi=2^{-n/2}\sum_x(-1)^{f(x)}\ket x\ket{g(x)},\]
其中$f$取0或1、g一一, 求第一端约化态及$\bra+^{\otimes n}\rho_A\ket+^{\otimes n}$; (2) $H^{\otimes n}R_aH^{\otimes n}$及$H^{\otimes n}T_aH^{\otimes n}$, $R_a\ket x=(-1)^{a\cdot x}\ket x,T_a\ket x=\ket{x\oplus a}$; (3) 素数p上$F_pR_aF_p^\dagger,F_pT_aF_p^\dagger$,$R_a\ket x=\e^{2\pi\ii ax/p}\ket x,T_a\ket x=\ket{x+1}$; (4) 原题声称$F_2\otimes F_3=PF_6P^{-1}$, P为置换矩阵; 检验并说明如何修正.\end{exercise}
\begin{solution}(1) 标签正交消交叉项,$\rho_A=I/2^n$, 因而$\bra+^{\otimes n}\rho_A\ket+^{\otimes n}=2^{-n}$. (2) 分别$T_a,R_a$. (3) $F_pR_aF_p^\dagger$使输出$y\mapsto y-a$, $F_pT_aF_p^\dagger=R_1$, 由指数求和或移指标直接得. (4) 原式若要求同一P, 并不成立: $\sqrt6 F_6$的对角元含$\e^{\ii\pi/3}$, 而$\sqrt6(F_2\otimes F_3)$的对角元集合为$\{1,\omega,\omega,-1,-\omega,-\omega\}$, $\omega=\e^{2\pi\ii/3}$, 置换相似不能改变对角元集合. 正确CRT式允许输入、输出两不同置换; 对模3端, 一个指标还须乘$2^{-1}=2\pmod3$, 即交换1、2. 用主作业二的角色分解逐项验即可得到正确的左右置换形式.
\end{solution}
\source{Optional PS6, 第1、3、4题, 第2题商业教材5.6.}
\begin{exercise}(1) 给$cn\log n$门近似n位QFT; (3) 四项一目标Grover一轮成功率; (4) $n=2^k$控制U, 用反极性Fredkin和单控制U的二叉递归, 给工作位和门数.\end{exercise}
\begin{solution}(1) 保留每位后的$b=\lceil\log_2(n/\epsilon)\rceil+O(1)$个受控小角旋转, 丢弃角和上界$O(n2^{-b})\le\epsilon$, 门$O(nb)$, 固定精度为$O(n\log n)$. (3) $\theta=\pi/6$, 一轮$\sin^2(3\theta)=1$. (4) 每对子控制a,b加工作0, 反极性Fredkin在a=0交换b和工作, 留原b位置值ab, 工作留$(1-a)b$. 因而可把各对子旗按二叉树合为全AND. n-1个工作, n-1个Fredkin, 一次由根控制U, 再逆树恢复所有输入和工作, 总$2n-1$门, 深度$2\log_2n+1$, 忽略每层可并行约定以外的连接开销.
\end{solution}

\source{Optional PS6第2题, 原页引用教材5.6并给出Fourier与相移实现加法的任务. 原句不逐字复刊, 以下按n位模$2^n$加法独立构造; 不把已有可读任务记成缺题面.}
\begin{exercise}给定已知整数a, 对n位计算基实现$\ket x\mapsto\ket{x+a\bmod2^n}$; 另给寄存器$\ket a$时, 实现可逆双寄存器加法并保留a. 使用QFT及相位门.\end{exercise}
\begin{solution}令$N=2^n$, $D_a\ket y=\e^{2\pi\ii ay/N}\ket y$. 则
\[F_N^\dagger D_aF_N\ket x=N^{-1}\sum_{z,y}\e^{2\pi\ii y(x+a-z)/N}\ket z=\ket{x+a\bmod N}.\]
已知a时,$D_a$按y二进制各位分解为单比特相位. 若a也为寄存器, 指数$ay/N$按两者各位的乘积展开, 使用两寄存器间受控相位, 并保留a; 对a叠加仍线性成立. 所有可能相位对最多$O(n^2)$, 加两次QFT亦为$O(n^2)$门. 若要普通整数和及进位, 先扩大足够位数避免模回绕; 本题选择的模加法条件已明确.\end{solution}

\originalsection{附加PS7--PS10: 模拟、回波与纠错}
\source{Optional PS7, 第1--5题.}
\begin{exercise}Hermitian A,B: (1) 验Trotter误差$O(1/n)$; (2) 从ZZ耦合和任意单位Pauli Hamiltonian生成所有两位Pauli乘积; (3) 生成任意多位Pauli串; (4) 用单轴旋转及$\exp(\ii\theta ZZ)$建CNOT并讨论最少两体耦合; (5) $\sigma_\pm=(X\pm\ii Y)/2$: 验$\cos(\omega t)X-\sin(\omega t)Y=\e^{\ii\omega t}\sigma_++\e^{-\ii\omega t}\sigma_-$及$\e^{\pm\ii\omega t}\sigma_\pm\e^{\ii\omega tZ/2}=\e^{\ii\omega tZ/2}\sigma_\pm$.\end{exercise}
\begin{solution}(1) 有界有限维时一步BCH误差$O(t^2\|[A,B]\|/n^2)$, n步至多$O(1/n)$; 无界算符需域条件, 不自动适用. (2) 单位旋转把Z换成X,Y, 或用$\ii[X_1, Z_1Z_2]\propto Y_1Z_2$, 再转第二位. (3) 对已有串$P$与只在一重叠位置反对易、并在新位非平凡的二体串Q取对易子, 得更长串; 归纳生成. (4) CZ等于$\e^{-\ii\pi/4}\e^{\ii\pi Z_1/4}\e^{\ii\pi Z_2/4}\e^{-\ii\pi Z_1Z_2/4}$, CNOT=$H_2$CZ$H_2$, 一次非局域耦合充分且至少一次必要. 官方解答把H看成绕$(x+z)/\sqrt2$的单次旋转, 给4个单比特门; 还可将目标端$H\e^{\ii\pi Z/4}$合成一个任意轴单比特旋转, 得
\[\mathrm{CNOT}_{1\to2}=\e^{-\ii\pi/4}\bigl(\e^{\ii\pi Z/4}\otimes H\e^{\ii\pi Z/4}\bigr)\e^{-\ii\pi Z_1Z_2/4}(I\otimes H).\]
在一次ZZ耦合、任意轴单比特旋转的约定下是3个单比特门. 两目标端都需基变换: 若只有前端或后端一个, 想令控制0分支为I会迫使该门对角, 控制1分支也只能对角, 无法是X. 控制端还须一次相位置换补偿两条件旋转的相对相位, 故这3个不能再省. 若只允许X、Y、Z轴旋转, 可用两$R_y(\pm\pi/2)$及两$R_z(\pi/2)$实现, 本卷给4个单比特门的上界, 不冒称该受限门库下的全局最小证明.
(5) $\e^{\ii\phi Z/2}\sigma_+\e^{-\ii\phi Z/2}=\e^{\ii\phi}\sigma_+$, 由$[Z,\sigma_+]=2\sigma_+$得; 把cos、sin驱动写升降项便得到原旋转系式.
\end{solution}
\source{Optional PS8, 第1--3题.}
\begin{exercise}(1) 写$U=\e^{-\ii\pi X_1/4}\e^{-\ii\pi Z_1Z_2/4}\e^{-\ii\pi Y_1/4}$的$4\times4$矩阵, 比较CNOT并补单比特旋转; (2) $H=(-\omega_1Z_1-\omega_2Z_2+JZ_1Z_2)/2$, 施第一位X的$\pi$脉冲回波两段, 证只演化第二位; (3) 三位所有Z及ZZ耦合, 设计解耦只留单比特项.\end{exercise}
\begin{solution}(1) 保留原题末门Y. 按$00,01,10,11$顺序, 先把三个指数分别写成$(I-\ii P)/\sqrt2$并依右到左相乘, 得
\[U=\frac1{\sqrt2}\begin{pmatrix}
1-\ii&0&0&0\\0&0&0&-1-\ii\\0&0&1+\ii&0\\0&1-\ii&0&0
\end{pmatrix}.\]
它与$\mathrm{CNOT}_{2\to1}$有同样的基态置换, 仅相位不同. 令$S=\operatorname{diag}(1,\ii)$, 按四列逐项比对得到
\[U=\e^{-\ii\pi/4}(S_1\otimes S_2^\dagger)\mathrm{CNOT}_{2\to1},\qquad
\mathrm{CNOT}_{2\to1}=\e^{\ii\pi/4}(S_1^\dagger\otimes S_2)U.\]
因此两输出端分别补$S^\dagger,S$即可, 整体相位无物理作用. 这里的下标仅表示门所在位, 张量中的S都是$2\times2$矩阵. 上述X--ZZ--Y序列已对照原页及官方答案, 数值QA验证此矩阵与补门.
(2) X1翻Z1及Z1Z2而不翻Z2, 各项对易, 两段各时长t相加只留$-\omega_2tZ_2$. 两个$\pi$脉冲自身给整体$-1$, 故$U=-\e^{\ii\omega_2tZ_2}$, 确只演化第二位. (3) 四段等时长的翻转符号选$+++,+--,-+-,--+$, 每对符号乘积总和0, 全ZZ抵消, 本地Z也总和0得到完整回波. 如要保留指定单位Z, 只平均另两位独立符号并保持指定一位正, 同样四段消相互作用和其他单位项. 给相邻符号差对应的X或Y $\pi$脉冲即可实现, 最后恢复原框架.
\end{solution}
\source{Optional PS9, 第1--6题.}
\begin{exercise}(1) 验$[x,p]=\ii\hbar$; (2) 原件取质量1, $H=(p^2+\omega^2x^2)/2,a=(\omega x+\ii p)/\sqrt{2\hbar\omega}$, 验$[a,a^\dagger]=1,[a,a^\dagger a]=a,[a^\dagger,a^\dagger a]=-a^\dagger,H=\hbar\omega(n+1/2)$; (3) $\hbar=\omega=1$时, 验$\psi_0(x)\propto\e^{-x^2/2}$满足$a\psi_0=0$; (4) $\ket n=(a^\dagger)^n\ket0/\sqrt{n!}$的升降矩阵元及能量; (5) JC模型$C=\omega(n+Z/2),D=g(a^\dagger\sigma_-+a\sigma_+)-(\Delta/2)Z$, 证$[C,D]=[C,H]=[D,H]=0$, $H=\hbar(C+D+\omega/2)$; (6) 四点方形图态.\end{exercise}
\begin{solution}(1) 对光滑测试函数用$p=-\ii\hbar\partial_x$, 乘法求导交叉项给$\ii\hbar$. (2--4) 为说明单位可推广到质量$m$, 取$m=1$即原题. $a=\sqrt{m\omega/(2\hbar)}x+\ii p/\sqrt{2m\hbar\omega}$, 得$[a,a^\dagger]=1$, $[n,a]=-a,[n,a^\dagger]=a^\dagger$. $a\psi_0=0$解$\psi_0=(m\omega/\pi\hbar)^{1/4}\e^{-m\omega x^2/(2\hbar)}$, 归一态$(a^\dagger)^n\ket0/\sqrt{n!}$, 升降系数$\sqrt{n+1},\sqrt n$, 能量$\hbar\omega(n+1/2)$. (5) 光子增加1伴Z/2减少1, 两作用总激发保持, 所以$[C,D]=0$, 其余由H线性分解得. (6) 边AB,BC,CD,DA,$\ket G=\frac14\sum_{a,b,c,d}(-1)^{ab+bc+cd+da}\ket{abcd}$, 直接CZ构造, 稳定子为$X_AZ_BZ_D$等四个循环生成元.
\end{solution}
\source{Optional PS10, 第1--4题.}
\begin{exercise}(1) 三位重复码可能落四个单X综合征子空间, 用两工作位建定位电路; (2) 偶三位串等和与奇三位串等和作为phase码字, 纠单Z; (3) Shor码: 每块0代表000、1代表111, $\ket{0_L}=(\ket{000}+\ket{011}+\ket{101}+\ket{110})_{\rm blocks}/2$, $\ket{1_L}=(\ket{001}+\ket{010}+\ket{100}+\ket{111})_{\rm blocks}/2$, 建纠任意单Pauli电路; (4) Hamming码, 令H为下述三行生成的二进制行空间, G为H与其全1陪集的并: $1111000,1100110,1010101$. 证逻辑负态编码为$\frac14(\sum_{x\in H}\ket x-\sum_{x\in G\setminus H}\ket x)$.\end{exercise}
\begin{solution}(1) CNOT1、2到辅助a, CNOT2、3到辅助b, 测a,b得到$e_1+e_2,e_2+e_3$并据对应X恢复. (2) 此码是$+++,$---的H变换后逻辑基等价, 测$X_1X_2,X_2X_3$可定位Z; 不测数据本身. (3) 每组三位先用Z对定位bit错并纠, 再测$X^{\otimes6}$覆盖块1+2、2+3定位phase所在块, 在该块任一位施Z. Y同时含bit、phase, 整体相位无关. 原件块标签偶、奇与常见Shor逻辑基相差逻辑H, 恢复的综合征结构相同. (4) CSS$\ket{0_L}=\sum_{x\in H}\ket x/\sqrt8$,$\ket{1_L}=\sum_{x\in H+\mathbf1}\ket x/\sqrt8$, 线性编码$\ket-=(0-1)/\sqrt2$即两陪集相减除$\sqrt{16}$. 本题原件明确给出的生成矩阵为$H=\bigl(\begin{smallmatrix}1&1&1&1&0&0&0\\1&1&0&0&1&1&0\\1&0&1&0&1&0&1\end{smallmatrix}\bigr)$. 其行空间有8项, 加全1得到互不交的8项陪集, 因每行重量为4而全1重量7. 所以上式16项均系数$\pm1/4$, 同时验证归一与原件的互补串描述.
\end{solution}

\originalsection{Quiz 1}
\source{Quiz 1, 第1--3题; 可参考公开Quiz 1 Solutions.}
\begin{exercise}(1) 绕$(x+z)/\sqrt2$轴$\pi$旋转与H比较; (2) $K=(XX+YY+ZZ+II)/2$, 求平方、$\e^{-\ii\pi K/4},\e^{-\ii\pi K/2}$、本征值及正交态; (3) $N>5$, $\psi=N^{-1/2}\sum_x\ket x\ket{3x\bmod N}\ket{5x\bmod N}$, 三端各QFT后等概率支持与概率.\end{exercise}
\begin{solution}(1) $R=-\ii(X+Z)/\sqrt2=-\ii H$, 只差整体相位. (2) K是SWAP,$K^2=I$, 两指数分别$(I-\ii K)/\sqrt2,-\ii K$. 对称00,11,$(01+10)/\sqrt2$值+1, 反对称单态值-1. (3) 输出$i,j,k$幅度$N^{-2}\sum_x\e^{2\pi\ii x(i+3j+5k)/N}$, 若$i+3j+5k=0\pmod N$为$1/N$, 否则0. 任意j,k唯一确定i, 共$N^2$项, 各概率$1/N^2$, 不要求3、5与N互素.
\end{solution}
\originalsection{Quiz 2}
\source{Quiz 2(test2), 第1--4题, 未取得独立官方答案.}
\begin{exercise}(1) 相位翻转$p$作用$\alpha0+\beta1$的态; (2) 四位码字$(0000+1111),(0110+1001),(1010+0101),(1100+0011)$各除$\sqrt2$, 找CSS的$C_1,C_2$, 是否检测一错; (3) 单X、单Z、单X再任意位置单Z是否都映到码正交空间; (4) 预共享Bell对受等概率三Pauli错误各$p/3$, 以$1-p$无错, 超密编码消息怎样传错?\end{exercise}
\begin{solution}(1) $\rho=\bigl(\begin{smallmatrix}|\alpha|^2&(1-2p)\alpha\beta^*\\(1-2p)\alpha^*\beta&|\beta|^2\end{smallmatrix}\bigr)$. (2) $C_1$为四位偶重量码$[4,3,2]$,$C_2=\{0000,1111\}$维1, 其对偶也是偶重量码, 故可检测单错, 编码2逻辑比特. 稳定子$XXXX,ZZZZ$. (3) 单X反对易ZZZZ, 单Z反对易XXXX; X后Z即使位置相同仍分别反对易两个生成元, 因而正交. 这不表示能纠正任意两错, 双X可以形成不被检测的逻辑算符. (4) 收到原消息概率$1-p$, 收到其余三个二比特标签各$p/3$, 即四元对称经典信道, 错误标签与本次Pauli超密编码叠加按模2相加.
\end{solution}
